\documentclass[10pt,letterpaper]{article}
\usepackage{cornell-notes}

\title{Clause 20.03.04.01: Class Template Out Ptr T}
\author{}
\date{}
\setDocTitle{Clause 20.03.04.01: Class Template Out Ptr T}
\setDocSubtitle{C++ 2024}
\setDocOwner{C++ 2024 Cornell Notes}
\setDocDate{\today}
\setCornellCollection{C++ 2024 Cornell Notes}
\setCornellUnitType{Clause}
\setCornellUnitNumber{20.03.04.01}
\setCornellUnitTitle{Class Template Out Ptr T}
\hypersetup{pdftitle={Clause 20.03.04.01: Class Template Out Ptr T},pdfsubject={Cornell notes},pdfauthor={}}

\begin{document}
\maketitle
\begin{CornellOverviewBox}{Overview}
\textbf{Classification:} Standard library - Normative\\[2pt]
\textbf{Learning objective:} Understand the rules, terminology, and semantic consequences associated with Class template out\_ptr\_t.
\end{CornellOverviewBox}\begin{CornellSummaryBox}{Summary}
{\sffamily\bfseries\color{Accent} Summary}\par\vspace{3pt}Section 20.3.4.1 focuses on Class template out\_ptr\_t. Apply the specified interface together with its mandates, effects, result conditions, exception behavior, and complexity constraints. When applying it, preserve the distinction between mandatory requirements, recommendations, permissions, and behavior deliberately left undefined, unspecified, or implementation-defined.
\end{CornellSummaryBox}

\end{document}
